\def\PrivateAuthorBlock{\begin{tabular}{c}
Weifeng Yang\\
% Yunnan Earthquake Agency, Kunming, China\\
\texttt{ywf841673182@gmail.com}

\end{tabular}}

\documentclass[11pt]{article}
\usepackage{amsmath,amssymb,amsthm}
\setlength{\textwidth}{6.3in}
\setlength{\oddsidemargin}{0.1in}
\setlength{\evensidemargin}{0.1in}
\setlength{\textheight}{8.6in}
\setlength{\topmargin}{-0.3in}
\setlength{\headheight}{14pt}
\setlength{\headsep}{22pt}
\setlength{\parindent}{1.2em}
\setlength{\parskip}{3pt}
\newcommand{\R}{\mathbb R}
\newcommand{\N}{\mathbb N}
\newcommand{\Nzero}{\mathbb N_0}
\newcommand{\gra}{\operatorname{gra}}
\newcommand{\dom}{\operatorname{dom}}
\newcommand{\sgn}{\operatorname{sgn}}
\newcommand{\pair}[2]{\langle #1,#2\rangle}
\newtheorem{theorem}{Theorem}
\newtheorem{proposition}{Proposition}
\newtheorem{lemma}{Lemma}
\newtheorem{example}{Example}
\newtheorem{corollary}{Corollary}
\renewcommand{\thecorollary}{C.\arabic{corollary}}
\newcommand{\dist}{\operatorname{dist}}
\pagestyle{plain}
\begin{document}

\begin{center}
{\Large\bfseries Nonmaximal sums of maximally monotone operators on a Banach space}\\[10pt]
\PrivateAuthorBlock\\[10pt]
% September 7, 2026\\[4pt]
 % Main manuscript
\end{center}


\begin{abstract}

We construct a maximally monotone operator on $c_0$ whose sum with a
bounded positive rank-one operator is not maximally monotone, although
the second operator has full domain and Rockafellar's interior-domain
condition holds. The construction couples countably many triangular
operators through a Lipschitz curve of block sums. Each actual parameter
has a finitely supported tail; the common limiting tail belongs to
$\ell^2\setminus\ell^1$. We compute the entire monotone polar and prove
maximality by showing that a point over the missing parameter interval
would require a nonsummable continuous-dual label. An exact pairing
identity then supplies a proper monotone extension of the sum. We also
prove an abstract criterion that identifies the precise failure of
endpoint realization used in the construction, and derive the finite
radial bound and the nonconvexity of the first operator's domain closure.
The same construction is realized on standard $\ell^1$ with its full
$\ell^\infty$ dual through a specified bounded surjection onto $c_0$.

\end{abstract}

\noindent\textbf{Keywords:} maximally monotone operator; sum theorem;
nonreflexive Banach space; monotone polar; rank-one operator.

\section{Introduction}
\label{sec:introduction}
Let $X$ be a real Banach space, let $X^*$ be its continuous dual, and let
$A,B:X\rightrightarrows X^*$ be maximally monotone operators. Their pointwise
sum is defined by $(A+B)x=\{a+b:a\in Ax,\ b\in Bx\}$. Rockafellar's sum
theorem \cite{rockafellar1970} states that $A+B$ is maximally monotone when
$X$ is reflexive and
\begin{equation}
\dom A\cap\operatorname{int}(\dom B)\ne\varnothing.
\label{eq:original-cq}
\end{equation}
The unrestricted sum question asks whether Eq.~\eqref{eq:original-cq}
alone is sufficient on an arbitrary real Banach space. The interior here
is taken in the norm topology of $X$, and maximality is taken in the full
pair $X\times X^*$.
The broader theory of maximality and convex representation on Banach
spaces is developed in Simons' monograph \cite{simons2008}.

Several positive results extend the theorem by imposing structure on one
of the factors. Yao \cite{yao2010} proves maximality when the first factor
is of type~(FPV) and the second factor has full domain. Borwein and Yao
\cite{borweinyao2012} prove the sum theorem under
Eq.~\eqref{eq:original-cq} when the first factor is a linear relation.
The factor order in the latter qualification matters: placing a
full-domain linear operator first would require interior points in the
domain of the other factor. These results motivate the question of
whether an elementary full-domain second factor can already produce a
nonmaximal sum with a nonlinear first factor.
Normal-cone partners under type~(FPV) hypotheses are treated by Voisei
\cite{voisei2010}. At the universal level, Bo\c{t}, Bueno, and Simons
\cite{botbuenosimons2016} show that restricting the sum question to pairs
with bounded common domain does not weaken it. This is a statement about
all admissible pairs, not a localization theorem for one fixed factor.

Convex representations provide a second way to organize these questions.
Fitzpatrick \cite{fitzpatrick1988} associates a convex function with a
monotone graph. Burachik and Svaiter \cite{burachiksvaiter2002} study the
associated family of convex representatives through enlargements, and
Penot \cite{penot2004} develops representation methods for compositions
and sums in the reflexive setting. For nonmaximal graphs,
representability and maximality must still be distinguished;
Mart\'{i}nez-Legaz and Svaiter \cite{martinezsvaiter2005} study this
distinction, including unique maximal extensions in finite dimensions.
Bueno, Mart\'{i}nez-Legaz,
and Svaiter \cite{buenomartinezsvaiter2014} compare representable and
monotone-polar closures. Our endpoint criterion concerns graphs containing
every point in specified affine fibres and computes their polar directly; it does not
replace these representation theories by an unrestricted normal form.
Marques Alves and Svaiter \cite{marquessvaiter2012} prove a sum theorem
under a qualification on domains of Fitzpatrick representations.
Such a qualification must be checked on those representations; it cannot
be replaced without proof by Eq.~\eqref{eq:original-cq}.

Nonreflexive constructions provide a concrete setting for the sum question.
Gossez's work \cite{gossez1971} develops monotone-operator theory beyond
reflexive spaces. The dense-type class has additional structure:
Marques Alves and Svaiter \cite{marquessvaiter2010} identify types~(D)
and~(NI), and Bauschke, Borwein, Wang, and Yao \cite{bauschkefp2012}
prove that Fitzpatrick--Phelps type~(FP) implies dense type.
Type~(FP) here is not type~(FPV). These results distinguish important
subclasses; none of these type assumptions is imposed in the question above.
Bueno and Svaiter \cite{buenosvaiter2011} construct a non-type~(D)
maximally monotone operator on $c_0$. Bauschke, Borwein, Wang, and Yao
\cite{bauschkeetal2011} develop a family of triangular operators and
further nonlinear examples. In particular, Example~4.1 of their
August~6, 2011 author version contains
the single-block triangular map used below. These earlier ingredients
give exact pairing identities. To obtain the sum behavior considered
here, we impose an additional constraint on the sums over countably many
blocks and verify maximality for every point of the full monotone polar.

We define this constraint through a $1$-Lipschitz map
$F:J\to\ell^1(\Nzero)\subset\ell^2(\Nzero)$, where
$J=(-\infty,-1)\cup(1,\infty)$. Its positive-index coordinates have finite
support for each actual parameter. At both deleted endpoints they tend,
coordinatewise and in $\ell^2$, to
$(1/(4n))_{n\ge1}$, which is not in $\ell^1$.
The triangular pairing identity first proves monotonicity. The complete
affine fibres of the graph then reduce the condition for an arbitrary
polar point to
\begin{equation}
\|F(t)-m_p\|_2^2+\nu^2\le(t-\tau)^2
\qquad(t\in J),
\label{eq:polar-intro}
\end{equation}
where $m_p\in\ell^1(\Nzero)$ is determined by its dual coordinate and
$\tau,\nu\in\R$. If $|\tau|>1$, the actual test $t=\tau$ returns a graph
point. If $|\tau|\le1$, limits in $\ell^2$ at the two endpoints
force the positive-index coordinates of $m_p$ to equal $1/(4n)$.
Their nonsummability excludes this second case and proves maximality.

The resulting first factor $A$ has nonempty domain at positive distance
from zero. Its construction also gives
$\langle x,a\rangle>-\sigma$ on $\gra A$, with
$\sigma=\pi^2/96<1/8$. We choose
$Px=\langle x,g\rangle g$ for a specified nonzero $g\in(c_0)^*$.
This factor is also the subdifferential of the continuous convex quadratic
$x\mapsto\langle x,g\rangle^2/2$, consistent with Rockafellar's
subdifferential maximality theorem \cite{rockafellar1970subdiff}.
We give a direct maximality proof for $P$ below.
The same pairing identity gives
$\langle x,a+Px\rangle>1-\sigma$ on every sum row.
Thus $(0,0)$ is monotonically related to the entire sum and is not in
its graph. The full domain of $P$ makes
Eq.~\eqref{eq:original-cq} immediate from nonemptiness of $\dom A$.

The construction on $c_0$ already addresses the unrestricted question.

We first prove an abstract criterion for maximality that computes the
entire monotone polar within the stated class of graphs. The criterion
explains why completion of the block-sum comparison space need not supply
an element of the continuous dual. We then verify its assumptions in the
coordinate construction, so no representation theorem for arbitrary
monotone operators is required.

We also transfer the construction to standard $\ell^1$ through a specified
bounded surjection $Q:\ell^1\to c_0$. Tests along $\ker Q$ force every
polar dual coordinate to factor through $Q$, and the remaining inequality
is exactly the one for $A$. We prove the transfer in the full pair
$\ell^1\times\ell^\infty$, using a controlled preimage for each point.

After the preliminaries in Section~\ref{sec:definitions},
Section~\ref{sec:endpoint-criterion} proves the exact endpoint criterion.
Section~\ref{sec:supporting} verifies its coordinate assumptions and proves
the supporting surjection results. Sections~\ref{sec:c0} and~\ref{sec:l1}
give the explicit $c_0$ example and its standard-$\ell^1$ realization.
Section~\ref{sec:domain-comparison} treats the domain closure and the
relevant published statements, and Section~\ref{sec:structure} explains
the roles of the affine fibres, the nonsummable limit and the second
factor. Further families and parameter variations are deferred from
this initial manuscript.


\section{Preliminaries}
\label{sec:definitions}
All scalars are real. Put $\N=\{1,2,\ldots\}$, $\Nzero=\{0,1,\ldots\}$,
and $I=\Nzero\times\N$. The space $Y=c_0(I)$ is global: for each
$\epsilon>0$, only finitely many coordinates have modulus at least
$\epsilon$. It has norm $\|\cdot\|_\infty$ and full dual $\ell^1(I)$,
with $\pair{x}{a}=\sum_{b,j}x_{b,j}a_{b,j}$.
The same pairing is used when the first argument lies in $\ell^\infty(I)$.
The unit coordinate vectors are $e_{b,j}$; the ambient space follows their use.

For a graph $G\subset X\times X^*$, define
\begin{equation}
G^\mu=\{(z,p)\in X\times X^*:
 \pair{z-x}{p-a}\ge0\quad\forall(x,a)\in G\}.
\tag{D1}\label{eq:D1}
\end{equation}
The graph is monotone when every two graph points have a nonnegative bracket,
and maximally monotone when it has no proper monotone extension in this same
full dual pair. For a monotone $G$, $G=G^\mu$ is equivalent to maximality.
Domain and range refer to actual graph points. Every inequality quantified
over the graph includes every value at a multivalued primal.


For the present graphs, which have no zero primal, define
\begin{equation}
\mathcal V(G)=\sup_{(x,a)\in G}
 \frac{\max\{0,-\pair{x}{a}\}}{\|x\|}.
\tag{D3}\label{eq:D3}
\end{equation}
At $(0,0)$ this is the nonnegative slope used by Verona and Verona
\cite{verona2008}. The finite upper bounds proved below are properties
of the constructed graphs, not an assumed general bound from that source.

We index coordinates by blocks. The map $m$ records the sum in each block,
and $E$ applies a triangular operator within that block. Its single-block
formula is the $\alpha=(1,1,\ldots)$ case of
\cite[author version, Example~4.1]{bauschkeetal2011}; the additional constraint below
couples the block sums through $F$.
For $a\in\ell^1(I)$, define the bounded maps
\begin{equation}
(ma)_b=\sum_{j\ge1}a_{b,j},\qquad
(Ea)_{b,j}=a_{b,j}+2\sum_{k>j}a_{b,k}.
\tag{D4}\label{eq:D4}
\end{equation}
Here $m:\ell^1(I)\to\ell^1(\Nzero)\subset\ell^2(\Nzero)$ and
$E:\ell^1(I)\to Y$. Put
\begin{equation}
\begin{aligned}
g&=e_{0,1}-e_{0,2}\in Y^*,&
h&=Eg=-e_{0,1}-e_{0,2}\in Y,\\
r(a)&=\pair{h}{a}=-a_{0,1}-a_{0,2},&
La&=-Ea+r(a)h.
\end{aligned}
\tag{D5}\label{eq:D5}
\end{equation}

For $s>0$ and $n\ge1$, write
\begin{equation}
\omega_s(n)=\frac{\max\{1-sn^{1/4},0\}}{4n},\quad
W(s)=\sum_{n\ge1}\omega_s(n)^2,\quad
\sigma=\sum_{n\ge1}\frac1{16n^2}=\frac{\pi^2}{96}.
\tag{D6}\label{eq:D6}
\end{equation}
Each $\omega_s$, $s>0$, has finite support. The symbol $\omega_0$ denotes only
the $\ell^2$ comparison vector $(1/(4n))_{n\ge1}$ and is never an actual
graph parameter. Define
\begin{equation}
\begin{aligned}
J&=(-\infty,-1)\cup(1,\infty),\\
F(t)&=(\sgn t,\omega_{|t|-1})\in\ell^1(\Nzero)\quad(t\in J),\\
D&=\{a\in\ell^1(I):r(a)\in J,\ ma=F(r(a))\}.
\end{aligned}
\tag{D7}\label{eq:D7}
\end{equation}

The operators and kernel are
\begin{equation}
\begin{aligned}
\gra A&=\{(La,a):a\in D\},\\
K&=\{k\in\ell^1(I):mk=0,\ r(k)=0\},\\
Px&=\pair{x}{g}g\quad(x\in Y).
\end{aligned}
\tag{D8}\label{eq:D8}
\end{equation}
For each $t\in J$, the following is one element of $D$, not the entire fibre:
\begin{equation}
a^t=-t e_{0,1}+(t+\sgn t)e_{0,3}
       +\sum_{n\ge1}\omega_{|t|-1}(n)e_{n,1}.
\tag{D9}\label{eq:D9}
\end{equation}

For standard $\ell^1$, let $Z=\ell^1(\N)$ and $Z^*=\ell^\infty(\N)$,
with their usual pairing. Fix once and for all an enumeration $(q_n)$ of
all finitely supported rational vectors in the closed unit ball of $Y$.
Define
\begin{equation}
\begin{aligned}
Qu&=\sum_{n\ge1}u_nq_n,& (Q^*a)_n&=\pair{q_n}{a},\\
\gra T&=\{(u,Q^*a):a\in D,\ Qu=La\},\\
f&=Q^*g,& Bu&=\pair{u}{f}f.
\end{aligned}
\tag{D10}\label{eq:D10}
\end{equation}
The proof establishes surjectivity with a norm-$2$ lift. No bounded linear
right inverse or complemented kernel is assumed. The polar for $T$ and
$T+B$ uses all $\ell^\infty$ labels.


\section{An exact endpoint criterion}
\label{sec:endpoint-criterion}
We first isolate the properties that determine the entire monotone polar.
The bilinear identity below makes a Lipschitz constraint monotone, and
the complete affine fibres reduce arbitrary polar points to a scalar
comparison. The two endpoint comparisons then identify exactly which
additional continuous-dual labels could extend the graph. The theorem
concerns the stated class of graphs; its assumptions will be verified
for the coordinate construction.

Let $X$ be a real Banach space and let $H_0$ be a real Hilbert space.
In this section put $H=\R\oplus H_0$. Let
$E:X^*\to X$ and $M:X^*\to H$ be bounded linear maps, and let
$g\in X^*\setminus\{0\}$ satisfy $Mg=0$ and
\begin{equation}
\pair{Ea}{b}+\pair{Eb}{a}=2(Ma,Mb)_H
\qquad(a,b\in X^*).
\label{eq:abs-bilinear}
\end{equation}
Define $h=Eg$, $r(a)=\pair{h}{a}$, $La=-Ea+r(a)h$, and
$K=\ker M\cap\ker r$. Assume the exact primal annihilator identity
\begin{equation}
\{x\in X:\pair{x}{k}=0\ \hbox{for all }k\in K\}=\R h.
\label{eq:abs-annihilator}
\end{equation}
Let $\omega:(0,\infty)\to H_0$ be $\ell$-Lipschitz, where
$0\le\ell\le1$, and suppose
\begin{equation}
\omega(s)\longrightarrow\eta\text{ in }H_0\quad(s\downarrow0),
\qquad 0<S:=\|\eta\|^2<1,\qquad \|\omega(s)\|^2\le S.
\label{eq:abs-tail}
\end{equation}
Put $J=(-\infty,-1)\cup(1,\infty)$ and
$F(t)=(\sgn t,\omega(|t|-1))$. For every $t\in J$, assume
there exists an actual $a^t\in X^*$ with $r(a^t)=t$ and
$Ma^t=F(t)$. We use every label with these values, defining
\begin{equation}
D=\{a\in X^*:r(a)\in J,\ Ma=F(r(a))\},\qquad
G=\{(La,a):a\in D\}.
\label{eq:abs-full-graph}
\end{equation}
In particular, each parameter fibre is $a^t+K$, without a boundedness
or continuity assumption on the choice $t\mapsto a^t$.

\renewcommand{\thetheorem}{E}
\begin{theorem}[Actual endpoint labels and maximality]
\label{thm:endpoint-criterion}
Under the preceding assumptions, define
\[
\widehat F(t)=
\begin{cases}
F(t),&|t|>1,\\
(t,\eta),&|t|\le1.
\end{cases}
\]
Then the entire monotone polar in $X\times X^*$ is
\begin{equation}
G^\mu=\{(Lp,p):p\in X^*,\ Mp=\widehat F(r(p))\}.
\label{eq:abs-entire-polar}
\end{equation}
It is the unique maximally monotone extension of $G$. Moreover, $G$
itself is maximally monotone if and only if
\begin{equation}
\nexists(p,\tau)\in X^*\times[-1,1]:
\qquad r(p)=\tau,\quad Mp=(\tau,\eta).
\label{eq:abs-endpoint-exclusion}
\end{equation}
When Eq.~\eqref{eq:abs-endpoint-exclusion} holds, the operator
$\mathcal A$ with graph $G$ omits zero and satisfies
$\mathcal V(G)\le S\|g\|$. With
$\mathcal P x=\pair{x}{g}g$, both $\mathcal A$ and $\mathcal P$
are maximally monotone, satisfy
$\dom\mathcal A\cap\operatorname{int}\dom\mathcal P\ne\varnothing$,
and have a nonmaximal sum: $(0,0)$ belongs to its monotone polar
and lies outside its graph.
\end{theorem}
\begin{proof}
Equation~\eqref{eq:abs-bilinear} at $a=b=g$ gives $r(g)=0$.
Using $Mg=0$ again gives
\begin{equation}
\pair{Ea}{g}=-r(a),\quad \pair{La}{g}=r(a),\quad
\pair{La-Lb}{a-b}=|r(a)-r(b)|^2-\|Ma-Mb\|^2.
\label{eq:abs-pairing}
\end{equation}
The actual label at $t=2$ also proves $h\ne0$.

Extend $\omega$ to zero by $\omega(0)=\eta$; its Lipschitz bound
persists by the norm limit. Define
$\chi(t)=\max\{-1,\min\{1,t\}\}$ and
$\zeta(t)=\max\{|t|-1,0\}$. Then
$\widehat F(t)=(\chi(t),\omega(\zeta(t)))$.
If $t,u$ lie on the same exterior branch, in the middle interval,
or one in the middle and one outside, the nonnegative changes in
$\chi$ and $\zeta$ add to $|t-u|$. On opposite exterior branches
they sum to $2+||t|-|u||\le|t-u|$. Thus in all cases
\[
|\chi(t)-\chi(u)|+|\zeta(t)-\zeta(u)|\le|t-u|,
\]
and
\[
\|\widehat F(t)-\widehat F(u)\|^2
\le|\chi(t)-\chi(u)|^2+\ell^2|\zeta(t)-\zeta(u)|^2
\le|t-u|^2.
\]
Equation~\eqref{eq:abs-pairing} proves monotonicity of both $G$
and the set on the right of Eq.~\eqref{eq:abs-entire-polar}.
It also proves that this latter set is contained in $G^\mu$.

For the reverse inclusion, take an arbitrary $(x,p)\in G^\mu$,
where $p$ belongs to the full continuous dual. For $a\in D$ and
$t=r(a)$, expansion using Eq.~\eqref{eq:abs-bilinear} gives
\begin{align*}
\pair{x-La}{p-a}
={}&\pair{x}{p}-\pair{x+Ep}{a}
       +2(Ma,Mp)-t r(p)+t^2-\|Ma\|^2.
\end{align*}
Replacing $a$ by $a+\theta k$, for all $\theta\in\R$ and
$k\in K$, forces $\pair{x+Ep}{k}=0$. Hence
$x=-Ep+v h$ by Eq.~\eqref{eq:abs-annihilator}. With
\[
\tau=\frac{v+r(p)}2,\qquad \nu=\frac{v-r(p)}2,
\]
the same expansion reduces exactly to
\begin{equation}
\pair{x-La}{p-a}=(t-\tau)^2-\nu^2-\|F(t)-Mp\|^2.
\label{eq:abs-polar-square}
\end{equation}
Every $t\in J$ is an actual test. If $|\tau|>1$, testing
$t=\tau$ forces $\nu=0$ and $Mp=F(\tau)$; thus
$v=r(p)=\tau$ and $x=Lp$.

If $|\tau|\le1$, write $Mp=(b,z)\in\R\oplus H_0$.
Limits along the two actual branches yield
\[
(b-1)^2+\|z-\eta\|^2+\nu^2\le(1-\tau)^2,
\qquad
(b+1)^2+\|z-\eta\|^2+\nu^2\le(1+\tau)^2.
\]
Multiply these by $(1+\tau)/2$ and $(1-\tau)/2$, respectively,
and add. Expansion gives
\[
(b-\tau)^2+\|z-\eta\|^2+\nu^2\le0.
\]
The endpoint weights remain valid at $\tau=\pm1$.
Consequently $Mp=(\tau,\eta)$, $r(p)=v=\tau$, and $x=Lp$.
This proves Eq.~\eqref{eq:abs-entire-polar}. The norm limits were
taken in $H$, not in $X^*$; the label $p$ was actual from the start.

We have proved that $G^\mu$ is monotone. Any point related to it
is related to $G$ and therefore belongs to $G^\mu$, so $G^\mu$
is maximally monotone. Every monotone extension of $G$ lies in
$G^\mu$; hence a maximal one must equal it. The missing part
$G^\mu\setminus G$ is exactly the set of labels in
Eq.~\eqref{eq:abs-endpoint-exclusion}. This proves the criterion.

Finally, each $a\in D$ gives
\[
|\pair{La}{g}|=|r(a)|>1,\quad
\|La\|>1/\|g\|,\quad
\pair{La}{a}=r(a)^2-1-\|\omega(|r(a)|-1)\|^2>-S.
\]
Thus zero is missing and, for every $a\in D$,
\[
\frac{\max\{0,-\pair{La}{a}\}}{\|La\|}
\le\frac{S}{\|La\|}<S\|g\|.
\]
Taking the supremum proves the stated whole-graph radial bound.
The map $\mathcal P$ is bounded, linear, positive, nonzero, and full-domain.
Its monotonicity follows from
$\pair{x-y}{\mathcal Px-\mathcal Py}=\pair{x-y}{g}^2$.
For an arbitrary point $(z,z^*)$ in its polar, tests at $z+tv$ give
\[
-t\pair{v}{z^*-\mathcal Pz}+t^2\pair{v}{g}^2\ge0
\qquad(t\in\R,\ v\in X).
\]
Both signs of arbitrarily small $t$ force $z^*=\mathcal Pz$.
Thus $\mathcal P$ is maximally monotone, and the actual label at
$t=2$ gives the ordered intersection condition. For every sum row,
\[
\pair{La}{a+\mathcal PLa}
=2r(a)^2-1-\|\omega(|r(a)|-1)\|^2>1-S>0.
\]
The sum is monotone and omits the zero primal. Adjoining $(0,0)$
is a proper monotone extension.
\end{proof}

The three conclusions use different parts of the argument. The Lipschitz
bound and Eq.~\eqref{eq:abs-pairing} prove monotonicity. The complete
fibres and Eq.~\eqref{eq:abs-endpoint-exclusion} prove maximality in the
original dual pair. Finally, the rank-one factor adds $r(a)^2$ to the
pairing, giving a positive bracket with the missing origin on every sum
row. A limit of profiles in $H$ contributes a polar point only when an
actual element of $X^*$ realizes both of its prescribed coordinates.

\section{Supporting lemmas}
\label{sec:supporting}
We now verify the assumptions of Theorem~\ref{thm:endpoint-criterion}
for the maps in Eqs.~\eqref{eq:D4}--\eqref{eq:D9}. The first three
lemmas supply the pairing identity, the exact primal annihilator and
the actual parameter fibres. We then prove the surjection facts used
to realize the same construction on standard $\ell^1$.

\begin{lemma}[The coordinate pairing identity]
\label{lem:coordinate-pairing}
The maps $m:\ell^1(I)\to\ell^1(\Nzero)\subset\ell^2(\Nzero)$ and
$E,L:\ell^1(I)\to c_0(I)$ are bounded. For all $a,c\in\ell^1(I)$,
Eqs.~\eqref{eq:C1} and~\eqref{eq:C2} below hold.
\end{lemma}
\begin{proof}
\label{sec:C1}
For $a\in\ell^1(I)$, $\|ma\|_1\le\|a\|_1$ and
$\|Ea\|_\infty\le2\|a\|_1$. If $a$ has finite support, then $Ea$ has
finite support: only finitely many blocks occur, and within a block no
output occurs after the last nonzero input coordinate. Approximation by
finite-support vectors and the bound prove $Ea\in c_0(I)$ globally.
Thus $L$ also maps $\ell^1(I)$ boundedly into $c_0(I)$.

For $a,c\in\ell^1(I)$, absolute convergence permits exchanging the double
sums. In each block the diagonal terms and the two triangular
off-diagonal sums give
\begin{equation}
\pair{Ea}{c}+\pair{Ec}{a}
 =2\sum_b\Big(\sum_j a_{b,j}\Big)\Big(\sum_j c_{b,j}\Big)
 =2\pair{ma}{mc}_{\ell^2}.
\tag{C1}\label{eq:C1}
\end{equation}
The absolute sum of all products is at most $\|a\|_1\|c\|_1$ before the
factor $2$, so this is an identity for all labels. Since $mg=0$, $h=Eg$,
and $r(g)=0$, Eq.~\eqref{eq:C1} implies $\pair{Ea}{g}=-r(a)$.
Consequently
\begin{equation}
\begin{aligned}
\pair{La}{a}&=r(a)^2-\|ma\|_2^2,&
\pair{La}{g}&=r(a),\\
\pair{La-Lc}{a-c}&=(r(a)-r(c))^2-\|ma-mc\|_2^2.
\end{aligned}
\tag{C2}\label{eq:C2}
\end{equation}
\end{proof}

\begin{lemma}[The exact primal annihilator]
\label{lem:coordinate-annihilator}
For $K$ in Eq.~\eqref{eq:D8},
\[
\{z\in c_0(I):\pair{z}{k}=0\ \text{for every }k\in K\}=\R h.
\]
\end{lemma}
\begin{proof}
Let $z\in c_0(I)$ globally annihilate $K$. In a block $b\ge1$, every
difference $e_{b,j}-e_{b,k}$ is in $K$. Thus $z$ is constant in that
entire block, and global $c_0$ forces that constant to vanish. In block
zero, differences between indices $j,k\ge3$ belong to $K$, so the common
tail value is also zero. Finally, $g=e_{0,1}-e_{0,2}$ belongs to $K$,
forcing $z_{0,1}=z_{0,2}$. It follows that $z$ is a scalar multiple of
$h$. Conversely, $\pair{h}{k}=r(k)=0$ proves that every multiple of $h$
annihilates $K$.
\end{proof}

\begin{lemma}[The tail and the complete parameter fibres]
\label{lem:coordinate-profile}
The map $s\mapsto\omega_s$ from $(0,\infty)$ to $\ell^2(\N)$
satisfies Eq.~\eqref{eq:abs-tail} with $\eta=\omega_0$, $S=\sigma<1/8$
and a Lipschitz constant strictly less than one. For every $t\in J$,
the label $a^t$ in Eq.~\eqref{eq:D9} belongs to $D$, and its entire
parameter fibre is $a^t+K$. No $p\in\ell^1(I)$ and $\tau\in[-1,1]$
satisfy $r(p)=\tau$ and $mp=(\tau,\omega_0)$.
\end{lemma}
\begin{proof}
\label{sec:C2}
For each $s>0$, $\omega_s(n)=0$ whenever $n\ge s^{-4}$. Also
$0\le\omega_s(n)\le1/(4n)$. Hence $W(s)\le\sigma<1/8$, and dominated
convergence in $\ell^2$ gives $\omega_s\to\omega_0$ and $W(s)\to\sigma$
as $s$ decreases to zero. The comparison vector $\omega_0$ is not in
$\ell^1$ by divergence of the harmonic series. No endpoint is added to $J$.

Since $v\mapsto\max\{v,0\}$ is $1$-Lipschitz on $\R$,
\begin{equation}
\|\omega_s-\omega_q\|_2^2\le c|s-q|^2,\qquad
c=\frac1{16}\sum_{n\ge1}n^{-3/2}<\frac3{16}<1.
\tag{C3}\label{eq:C3}
\end{equation}
The strict numerical bound follows from
$\sum_{n\ge1}n^{-3/2}<1+\int_1^\infty x^{-3/2}\,dx=3$.
For $t,u$ on the same branch of $J$, Eq.~\eqref{eq:C3} gives
$\|F(t)-F(u)\|_2\le|t-u|$. For $t=1+s$ and $u=-1-q$, the squared
distance is at most $4+c(s-q)^2\le(2+s+q)^2$, and the opposite
orientation is identical. Thus $F$ is $1$-Lipschitz on all of $J$.

The finitely supported vector $a^t$ in Eq.~\eqref{eq:D9} satisfies $r(a^t)=t$ and
$ma^t=F(t)$, proving nonemptiness for every actual $t$. The complete
fibre $\{a:r(a)=t,\ ma=F(t)\}$ is precisely $a^t+K$.
For every $p\in\ell^1(I)$, the vector $mp$ is in $\ell^1(\Nzero)$.
Since $\omega_0\notin\ell^1(\N)$, the last assertion follows.
\end{proof}

The following pullback argument is used only to transport the construction.
We include its proof to specify the role of arbitrary continuous-dual
labels and of the controlled preimages.

\begin{lemma}[Pullback by a bounded surjection]
\label{lem:surjection-transport}
Let $U,V$ be real Banach spaces and let $Q:U\to V$ be a bounded
linear surjection. Suppose $C_Q>0$ and every $y\in V$ has a preimage
$u\in U$ with $\|u\|\le C_Q\|y\|$. If
$\mathcal M:V\rightrightarrows V^*$ is maximally monotone, then
\[
\gra(Q^*\mathcal M Q)
=\{(u,Q^*a):(Qu,a)\in\gra\mathcal M\}
\]
is maximally monotone in the full pair $U\times U^*$.
\end{lemma}
\begin{proof}
A maximally monotone graph is nonempty, since the empty graph admits
a one-point monotone extension. Surjectivity therefore makes the
pullback graph nonempty. For two of its rows,
\[
\pair{u-v}{Q^*(a-b)}=\pair{Qu-Qv}{a-b}\ge0,
\]
so it is monotone.

Let $(v,v^*)\in U\times U^*$ be related to this entire graph. Fix one
row $(u,Q^*a)$. For each $k\in\ker Q$, all $(u+t k,Q^*a)$ are rows,
and their polar inequalities have the form
\[
\pair{v-u}{v^*-Q^*a}-t\pair{k}{v^*}\ge0
\qquad(t\in\R).
\]
Thus $v^*$ annihilates $\ker Q$. Define $p(y)=\pair{u}{v^*}$ when
$Qu=y$. Differences of preimages prove well-definedness, and linear
combinations of preimages prove linearity. The prescribed lifting bound
gives $|p(y)|\le C_Q\|v^*\|\|y\|$, so $p\in V^*$ and $v^*=Q^*p$.
Every $(y,a)\in\gra\mathcal M$ has a preimage row. The remaining
inequalities are
\[
\pair{Qv-y}{p-a}\ge0\qquad((y,a)\in\gra\mathcal M).
\]
Maximality of $\mathcal M$ gives $(Qv,p)\in\gra\mathcal M$, and hence
$(v,v^*)$ belongs to the pullback graph.
\end{proof}

\begin{lemma}[The specified quotient onto $c_0$]
\label{lem:specified-quotient}
The map $Q:Z=\ell^1(\N)\to Y=c_0(I)$ of Eq.~\eqref{eq:D10} is a
bounded surjection with $\|Q\|=1$. Its adjoint satisfies
$\|Q^*a\|_\infty=\|a\|_1$ for every $a\in\ell^1(I)$.
For the fixed enumeration $(q_n)$, a map $\mathcal R:Y\to Z$ can be
specified with
\[
Q\mathcal R(x)=x,\qquad \|\mathcal R(x)\|_1\le2\|x\|_\infty.
\]
\end{lemma}
\begin{proof}
\label{sec:L1}
Finitely supported rational vectors in the closed unit ball of $Y$ are
countable and dense in that ball: truncate a $c_0$ vector and approximate
each retained coordinate by a rational inside $[-1,1]$. Fix an enumeration
of all these vectors as $(q_n)$. For $u\in Z=\ell^1(\N)$, the sum defining
$Qu$ converges absolutely in the Banach space $Y$, and
$\|Qu\|_\infty\le\|u\|_1$. Coordinate unit vectors occur among the $q_n$,
so $\|Q\|=1$.

For an arbitrary $x\in Y$, set $r_0=x$. If $r_k\ne0$, choose the least index $n_k$ for which the actual
$q_{n_k}$ satisfies $\|q_{n_k}-r_k/\|r_k\|_\infty\|_\infty<1/2$ and set
\begin{equation}
c_k=\|r_k\|_\infty,\qquad r_{k+1}=r_k-c_kq_{n_k}.
\tag{L1}\label{eq:L1}
\end{equation}
If a residual vanishes, stop. Otherwise
$\|r_k\|_\infty\le2^{-k}\|x\|_\infty$ and $\sum_k c_k\le2\|x\|_\infty$.
Define $u_n=\sum_{k:n_k=n}c_k$. Repeated indices are combined.
Positivity of the $c_k$ and summability give $u\in\ell^1$ with
$\|u\|_1=\sum_k c_k\le2\|x\|_\infty$.
Telescoping Eq.~\eqref{eq:L1}, then absolute convergence under regrouping,
proves
\begin{equation}
Qu=x,\qquad \|u\|_1\le2\|x\|_\infty.
\tag{L2}\label{eq:L2}
\end{equation}
For $x=0$, use $u=0$. This proves surjectivity with the asserted lifting
bound. The construction does not assert that $x\mapsto u$ is linear or
continuous, nor that $\ker Q$ is complemented.

\label{sec:L2}
Absolute convergence gives $\pair{u}{Q^*a}=\pair{Qu}{a}$ and
$\|Q^*a\|_\infty\le\|a\|_1$. For each finite subset of $I$, its sign
vector for $a$ is a finitely supported rational vector of norm at most
one and occurs in the enumeration. Taking these finite sign tests proves
\begin{equation}
\|Q^*a\|_\infty=\|a\|_1\qquad(a\in\ell^1(I)).
\tag{L3}\label{eq:L3}
\end{equation}
In particular, $Q^*$ is injective. The least-index choices in
Eq.~\eqref{eq:L1} specify
\[
\mathcal R(x)=\sum_k c_k e_{n_k},\qquad \mathcal R(0)=0.
\]
The series converges in $\ell^1$ by the bound on $\sum_k c_k$,
and Eq.~\eqref{eq:L2} gives both asserted properties.
\end{proof}

For this fixed enumeration,
\[
\gra T=\{(\mathcal R(La)+v,Q^*a):a\in D,\ v\in\ker Q\}.
\]
Indeed, $Q\mathcal R(La)=La$, and any other preimage differs by
an element of $\ker Q$. Thus the entire graph is specified by
Eq.~\eqref{eq:D10}; its definition does not select a maximal extension.
This description asserts neither a bounded linear or continuous right
inverse nor decidability of membership in the infinite-dimensional graph.

\section{A counterexample on $c_0$}
\label{sec:c0}
\label{sec:statements}
\renewcommand{\theexample}{C}
\begin{example}[$c_0$ realization]\label{thm:C}
Let $I=\Nzero\times\N$ and $Y=c_0(I)$ with its usual supremum norm and
full dual $Y^*=\ell^1(I)$. Let $A:Y\rightrightarrows Y^*$ and the bounded
positive rank-one map $P:Y\to Y^*$ be defined by Eq.~\eqref{eq:D8}.
Then:
\begin{enumerate}
\item $A$ and $P$ are maximally monotone in $Y\times Y^*$.
\item $\dom P=Y$ and $\dom A\ne\varnothing$, so
$\dom A\cap\operatorname{int}(\dom P)\ne\varnothing$.
\item Every $(x,a)\in\gra A$ satisfies $\|x\|_\infty>1/2$ and
$\pair{x}{a}\ge-2\sigma\|x\|_\infty$, where $\sigma=\pi^2/96<1/8$.
\item Every $(x,b)\in\gra(A+P)$ satisfies
$\pair{x}{b}>1-\sigma>7/8$.
\item $(0,0)$ is monotonically related to the entire graph of $A+P$ and
does not belong to that graph. Hence $A+P$ is not maximally monotone.
\end{enumerate}
\end{example}
\begin{proof}
Lemmas~\ref{lem:coordinate-pairing}--\ref{lem:coordinate-profile}
verify every assumption of Theorem~\ref{thm:endpoint-criterion}
with $X=Y$, $M=m$, $H_0=\ell^2(\N)$, $\omega(s)=\omega_s$,
$\eta=\omega_0$ and $S=\sigma$. In particular, $g\ne0$ and $mg=0$
by Eq.~\eqref{eq:D5}; all parameter fibres are present and the
exact primal annihilator is $\R h$.

By Eq.~\eqref{eq:C2}, the Lipschitz constraint proves monotonicity
for every pair of labels in $D$. The full polar formula
\eqref{eq:abs-entire-polar} and the nonsummability of $\omega_0$
then give $(\gra A)^\mu=\gra A$: no actual $p\in\ell^1(I)$ can
satisfy $mp=(\tau,\omega_0)$ for $|\tau|\le1$. This proves
maximality in $c_0(I)\times\ell^1(I)$.

\label{sec:C5}
For every $a\in D$, $t=r(a)$ satisfies $|t|>1$. By Eq.~\eqref{eq:C2}
and $\|g\|_1=2$,
\begin{equation}
\|La\|_\infty\ge|t|/2>1/2,\qquad
\pair{La}{a}=t^2-1-W(|t|-1)>-\sigma.
\tag{C9}\label{eq:C9}
\end{equation}
It follows that $A(0)=\varnothing$ and
$\pair{La}{a}\ge-2\sigma\|La\|_\infty$ for every actual graph label.
In particular $\mathcal V(\gra A)\le2\sigma<1/4$.
Indeed, for each $a\in D$, Eq.~\eqref{eq:C9} gives
$\max\{0,-\pair{La}{a}\}/\|La\|_\infty
\le\sigma/\|La\|_\infty<2\sigma$; taking the supremum proves the bound.

\label{sec:C6}
The rank-one assertion and maximality of $P$ follow from the direct
proof in Theorem~\ref{thm:endpoint-criterion}. Here $\|g\|_1=2$,
so $\|P\|\le4$, and $Pe_{0,1}=g\ne0$.
The actual label $a^2=-2e_{0,1}+3e_{0,3}$ gives
$La^2=-6e_{0,1}-8e_{0,2}-3e_{0,3}$. Thus
$\dom A\cap\operatorname{int}(\dom P)\ne\varnothing$, with
$\operatorname{int}(\dom P)=Y$. Every sum point has the form
$(La,a+tg)$, $a\in D$ and $t=r(a)$, and
\begin{equation}
\pair{-La}{-a-tg}=2t^2-1-W(|t|-1)>1-\sigma>7/8.
\tag{C11}\label{eq:C11}
\end{equation}
This proves that $(0,0)$ is related to every sum point. The sum is
monotone as the sum of two monotone operators; it has no zero primal by
Eq.~\eqref{eq:C9}. Adding $(0,0)$ is therefore a proper monotone extension.

\end{proof}

\section{The counterexample on standard $\ell^1$}
\label{sec:l1}
\renewcommand{\theexample}{L}
\begin{example}[Standard $\ell^1$ realization]\label{thm:L}
On $Z=\ell^1(\N)$ with its usual norm and full dual $Z^*=\ell^\infty(\N)$,
let $T:Z\rightrightarrows Z^*$ and $B:Z\to Z^*$ be defined by Eq.~\eqref{eq:D10}.
Then:
\begin{enumerate}
\item $T$ is maximally monotone in the full pair $Z\times Z^*$.
\item $B$ is nonzero, bounded, positive, rank-one, and maximally monotone,
with $\dom B=Z$. In particular,
$\dom T\cap\operatorname{int}(\dom B)\ne\varnothing$.
\item $T(0)=\varnothing$, and every $(u,b)\in\gra T$ satisfies
$\|u\|_1>1/2$ and $\pair{u}{b}\ge-2\sigma\|u\|_1$.
Thus $\mathcal V(\gra T)\le2\sigma<1/4$.
\item Every $(u,c)\in\gra(T+B)$ satisfies
$\pair{u}{c}>1-\sigma>7/8$.
\item $(0,0)$ belongs to the full monotone polar of $\gra(T+B)$ and does
not belong to $\gra(T+B)$. Hence $T+B$ is not maximally monotone.
\end{enumerate}
\end{example}
\begin{proof}
Lemma~\ref{lem:specified-quotient} gives a surjection $Q$ with the
norm-$2$ lifting bound, and Example~\ref{thm:C} proves maximality
of $A$. Thus Lemma~\ref{lem:surjection-transport} applies to $T$.
For the specified spaces, we record the full-dual reduction explicitly.

The graph of $T$ is nonempty by
Eq.~\eqref{eq:L2} and nonemptiness of $D$. It is monotone because the
bracket of any two rows $(u,Q^*a)$ and $(v,Q^*c)$ equals
$\pair{La-Lc}{a-c}$, which is nonnegative by Eq.~\eqref{eq:C2} and the Lipschitz bound.

Let $(v,v^*)$ be a polar point with $v^*$ an arbitrary element of
$\ell^\infty(\N)$. Each actual fibre of $T$ contains $u+\lambda k$ for
every $k\in\ker Q$ and every $\lambda\in\R$. Testing those rows forces
$\pair{k}{v^*}=0$ for every such $k$. Define a functional on $Y$ by
\begin{equation}
\lambda(x)=\pair{u}{v^*}\quad\text{when }Qu=x.
\tag{L4}\label{eq:L4}
\end{equation}
Two preimages differ by $\ker Q$, proving well-definedness. Linear
combinations of preimages prove linearity. For each $x$, choose the
preimage in Eq.~\eqref{eq:L2}; then
$|\lambda(x)|\le2\|v^*\|_\infty\|x\|_\infty$. Thus $\lambda$ belongs to
the continuous dual of $c_0(I)$.

Explicitly put $p_i=\lambda(e_i)$. Finite sign tests give
$\sum_{i\in H}|p_i|\le\|\lambda\|$ for every finite $H$, whence
$p\in\ell^1(I)$. Equality $\lambda(x)=\pair{x}{p}$ holds first for
finite-support $x$ and then for every $c_0$ vector by density and continuity.
Therefore $v^*=Q^*p$: the two functionals agree on every $u$ because
Eq.~\eqref{eq:L4} applies to $Qu$.

Every original $a\in D$ admits a lift $u$ of $La$ by Eq.~\eqref{eq:L2}.
The full $T$-polar inequality now reads
\begin{equation}
\pair{Qv-La}{p-a}\ge0\qquad\forall a\in D.
\tag{L5}\label{eq:L5}
\end{equation}
Example~\ref{thm:C} gives $p\in D$ and $Qv=Lp$. Hence
$(v,v^*)=(v,Q^*p)$ is an actual $T$ row. This proves maximality in the
full standard pair $\ell^1\times\ell^\infty$. Membership of $v^*$ in
$\operatorname{range}(Q^*)$ was derived from the actual kernel tests.

\noindent\textit{The radial bound on every graph value.}
\label{sec:L3}
For every $(u,Q^*a)\in\gra T$, with $t=r(a)$,
\begin{equation}
\begin{aligned}
\pair{u}{Q^*a}&=\pair{La}{a}
 =t^2-1-W(|t|-1)>-\sigma,\\
\|u\|_1&\ge\|Qu\|_\infty=\|La\|_\infty>1/2.
\end{aligned}
\tag{L6}\label{eq:L6}
\end{equation}
Thus $T(0)=\varnothing$ and $\pair{u}{Q^*a}\ge-2\sigma\|u\|_1$ on
the entire graph, including every $\ker Q$ translate. Consequently,
$\mathcal V(\gra T)\le2\sigma<1/4$. No exact equality for this
norm-dependent supremum is asserted.

\noindent\textit{The rank-one factor and the intersection condition.}
\label{sec:L4}
Let $f=Q^*g$. By Eq.~\eqref{eq:L3}, $\|f\|_\infty=\|g\|_1=2$, so $f$
is nonzero. The map $B:u\mapsto\pair{u}{f}f$ is nonzero, linear, bounded
with norm at most $4$, positive, and rank-one. The rank-one proof in
Theorem~\ref{thm:endpoint-criterion}, now testing all $v\in\ell^1$
and labels in $\ell^\infty$, proves $B$ maximally monotone.
Its domain is all of $Z$, so $\operatorname{int}(\dom B)=Z$.

The vector $q_\circ=-(3/4)e_{0,1}-e_{0,2}-(3/8)e_{0,3}$ occurs as
$q_{n_0}$. For $u_\circ=8e_{n_0}$, $Qu_\circ=La^2$.
Thus $(u_\circ,Q^*a^2)$ is an actual $T$ row, and
$\dom T\cap\operatorname{int}(\dom B)\ne\varnothing$ in the stated
factor order.

\noindent\textit{A strict polar witness outside the graph.}
\label{sec:L5}
For every row $(u,Q^*a)$ of $T$, let $t=r(a)$. Then
$\pair{u}{f}=\pair{Qu}{g}=\pair{La}{g}=t$.
Consequently each sum row has work
\begin{equation}
\pair{u}{Q^*a+Bu}
 =2t^2-1-W(|t|-1)>1-\sigma>7/8.
\tag{L7}\label{eq:L7}
\end{equation}
The sum is monotone by monotonicity of its factors.
Equation~\eqref{eq:L7} shows that $(0,0)$ is related to all its rows.
Since $T(0)=\varnothing$, $(0,0)$ is not a sum row. Adjoining this point
gives a proper monotone extension, completing the proof of
Example~\ref{thm:L} with the standard Banach space, full dual, and ordered
intersection condition. No theorem about sums was used to prove either
factor maximal.
\end{proof}

The specified map $Q$ preserves the pairing and transports the same
missing polar point. The two examples realize one construction on
different Banach spaces, with maximality proved in each full dual pair.

\section{Domain geometry and an almost-convexity statement}
\label{sec:domain-comparison}
The same graph also has a nonconvex domain closure. This consequence
helps identify the relationship between the construction and claims
about the domains of arbitrary maximally monotone operators.

\begin{proposition}\label{prop:domain-gap}
For the operator $A$ of Example~\ref{thm:C},
$0\in\operatorname{conv}(\dom A)$ and
$\operatorname{dist}(0,\dom A)\ge1/2$.
In particular, $\overline{\dom A}$ is not convex.
\end{proposition}
\begin{proof}
At $t=2$ and $t=-2$, the tails in Eq.~\eqref{eq:D9} vanish, so that
$a^{-2}=-a^2$. Since $L$ is linear, both $La^2$ and $-La^2$ belong
to $\dom A$. Their midpoint is zero. Equation~\eqref{eq:C9} gives
$\|x\|_\infty>1/2$ for every $x\in\dom A$, and hence the asserted
distance bound. A convex closure containing $La^2$ and $-La^2$ would
contain zero, contradicting this bound.
\end{proof}

Eberhard and Wenczel \cite[Theorem~57]{eberhardwenczel2024} state that
the norm closure of the domain of every maximally monotone operator on
a Banach space is convex. Proposition~\ref{prop:domain-gap}, together
with the proof of Example~\ref{thm:C}, gives an incompatible
conclusion for the displayed operator. The proof on
\cite[pp.~39--40]{eberhardwenczel2024} uses a small-parameter estimate
from Proposition~56, Eq.~(65), along a family in which the operator and
the evaluated point vary. The estimate for a fixed operator and point
has an object-dependent threshold; it does not supply a threshold
uniform over that family. The additional uniform estimate is needed
for that application. Our construction and maximality argument do not
use this published domain-closure statement or its negation.

There is also a version-specific issue with the separable-Asplund sum
statement in Bo\c{t} and Csetnek~\cite[Corollary~11]{botcsetnek2008}.
Its written hypotheses apply to the $c_0(I)$ construction: this space
is separable Asplund and $\dom P=Y$, so the convexified domain difference
is all of $Y$. The inspected proof invokes their Theorem~3, which asserts
bidual conjugate domination for representatives of every maximally
monotone operator when the dual space is separable. The operator in
\cite[author version, Lemma~2.1]{buenosvaiter2011} shows why that assertion cannot be used.
Write $G$ for the Gossez skew map on $\ell^1$, $e=(1,1,\ldots)$, and
$s(a)=\sum_n a_n$. That operator has graph
$\{(-Ga,a):a\in\ell^1,\ s(a)=0\}$. Its Fitzpatrick function is the
indicator of $x+Gb\in\R e$; since $x\in c_0$ and $Gb$ tends to $-s(b)$,
the finite-value points satisfy $x=-Gb-s(b)e$. For its conjugate on
$\ell^1\times\ell^\infty$, direct substitution therefore gives
\begin{equation}
\begin{aligned}
F_S^{c*}(p,\xi)
 &=\sup_{b\in\ell^1}\pair{Gp-s(p)e+\xi}{b},\\
F_S^{c*}(e_1,e-Ge_1)&=0<1=\pair{e-Ge_1}{e_1}.
\end{aligned}
\label{eq:bidual-interface}
\end{equation}
This calculation refutes the domination assertion used in that proof;
failure of the proof alone does not decide the standalone sum statement.
Example~\ref{thm:C} is proved independently of it.

The standard-$\ell^1$ realization in Example~\ref{thm:L} uses
the maximality of $A$ established in Example~\ref{thm:C}; the change
of space preserves that dependency. The distinction concerns the cited
2008 preprint, not an assertion about every later version or possible repair.


\section{The mechanism of the counterexample}
\label{sec:structure}
The construction uses the coordinates $r(a)$ and $ma$ to express the
symmetric part of the graph pairing:
\[
\pair{La-Lb}{a-b}=|r(a)-r(b)|^2-\|ma-mb\|_2^2.
\]
A Lipschitz curve in these coordinates makes the graph monotone.
Maximality requires the additional information supplied by the complete
affine fibres. Testing every kernel direction forces an arbitrary polar
query to have the scalar form in Eq.~\eqref{eq:abs-polar-square}.
At an exterior parameter this form returns an actual graph point.
Over the middle interval, the two endpoint comparisons force the
completed profile $(\tau,\eta)$.

Theorem~\ref{thm:endpoint-criterion} gives the exact distinction:
the completed profile contributes a polar point precisely when
some $p\in X^*$ satisfies both $r(p)=\tau$ and $Mp=(\tau,\eta)$.
For the constructed operator, $mp\in\ell^1$ and
$\eta=(1/(4n))_{n\ge1}\notin\ell^1$. Thus the comparison geometry
has a middle segment with no continuous-dual labels realizing it.
The proof takes limits in the comparison space; it never enlarges
the dual space of $c_0$.

This failure of realization can also be seen along actual graph rows.
Set $t_n^\pm=\pm(1+1/n)$, $a_n^\pm=a^{t_n^\pm}$, and
$x_n^\pm=La_n^\pm$. For $t\in J$ and $\epsilon=\sgn t$,
direct substitution in Eqs.~\eqref{eq:D4}--\eqref{eq:D9} gives
\[
(La^t)_{0,1}=-2t-2\epsilon,\qquad
(La^t)_{0,2}=-3t-2\epsilon,\qquad
(La^t)_{0,3}=-t-\epsilon,
\]
with the remaining block-zero coordinates zero. In block $b\ge1$,
the only nonzero coordinate is $-\omega_{|t|-1}(b)$ at index $1$.
Consequently, $x_n^\pm$ converge in the supremum norm to
$x_\pm\in c_0(I)$ with block-zero coordinates $\mp(4,5,2)$ and
positive-block first coordinates $-1/(4b)$. Indeed, convergence of
the tails in $\ell^2$ implies convergence in the supremum norm.
The limits satisfy $\pair{x_\pm}{g}=\pm1$ and lie outside $\dom A$.
On the other hand,
\[
\|a_n^\pm\|_1=2(1+1/n)+1+\|\omega_{1/n}\|_1\longrightarrow\infty.
\]
For each fixed $N$, the sum of the first $N$ tail coordinates tends to
$\sum_{b=1}^N1/(4b)$; these partial sums are unbounded. This proves
the divergence without asserting convergence of the labels in the dual.

The sum obstruction uses a different consequence of the pairing identity.
Every row of the first graph satisfies
\[
\pair{La}{a}=r(a)^2-1-W(|r(a)|-1)>-\sigma,\qquad
\|La\|_\infty>1/2.
\]
These inequalities prove $\mathcal V(\gra A)\le2\sigma$ directly.
The second factor adds exactly $r(a)^2$, making the pairing positive
on the entire sum graph. Since zero remains outside the domain,
$(0,0)$ is a missing polar point of the sum. The finite radial bound
is thus an output of the construction, not an assumed universal
property or an unproved consequence of the desired counterexample.


\section{Conclusions}
\label{sec:conclusion}

We constructed a maximally monotone operator on $c_0$ whose sum with a
full-domain positive rank-one operator is not maximally monotone under
the classical interior-domain condition. The proof combines an exact
triangular pairing identity, complete affine fibres and two endpoint
comparisons that would force an unavailable continuous-dual label.
The same identity proves the finite radial bound and verifies the
missing polar point of the sum directly. The abstract criterion
identifies exactly when the endpoint comparison adds a graph point,
separating the Hilbert-space comparison from realization in the
original dual pair. A specified bounded surjection transports this
construction and its missing sum point to standard $\ell^1$, with
maximality proved in the full $\ell^1\times\ell^\infty$ pair.


\section*{Use of AI}
The author directed this research, fixed the unrestricted problem and
admissible assumptions, and led the development of its tools and structural
framework by setting mathematical requirements, challenging incomplete
arguments, and selecting or rejecting proposed mechanisms. In particular,
the author required separate explanations of monotonicity, maximality and
sum failure, proofs in the full continuous duals, and explicit $c_0$ and
standard $\ell^1$ realizations. Under this direction, ChatGPT and Codex
assisted with candidate constructions, proof development and reconstruction,
literature checks, adversarial review, Lean formalization of the $c_0$
counterexample, and LaTeX exposition. The author determined the presentation
and the scope of the results included here and retains responsibility for
the manuscript and all submission decisions.

\begin{thebibliography}{99}
\raggedright
\bibitem{simons2008}
S. Simons,
\emph{From Hahn--Banach to Monotonicity}, 2nd ed.,
Lecture Notes in Mathematics, vol.~1693, Springer, 2008.
doi:10.1007/978-1-4020-6919-2.

\bibitem{fitzpatrick1988}
S. Fitzpatrick,
Representing monotone operators by convex functions,
in: Workshop/Miniconference on Functional Analysis and Optimization,
Proceedings of the Centre for Mathematical Analysis, Australian National
University, vol.~20, 1988, pp.~59--65.

\bibitem{burachiksvaiter2002}
R.~S. Burachik and B.~F. Svaiter,
Maximal monotone operators, convex functions and a special family of enlargements,
\emph{Set-Valued Analysis} \textbf{10} (2002), 297--316.
doi:10.1023/A:1020639314056.

\bibitem{penot2004}
J.-P. Penot,
The relevance of convex analysis for the study of monotonicity,
\emph{Nonlinear Analysis: Theory, Methods \& Applications}
\textbf{58} (2004), 855--871.
doi:10.1016/j.na.2004.05.018.

\bibitem{martinezsvaiter2005}
J.-E. Mart\'{i}nez-Legaz and B.~F. Svaiter,
Monotone operators representable by l.s.c. convex functions,
\emph{Set-Valued Analysis} \textbf{13} (2005), 21--46.
doi:10.1007/s11228-004-4170-4.

\bibitem{buenomartinezsvaiter2014}
O. Bueno, J.-E. Mart\'{i}nez-Legaz, and B.~F. Svaiter,
On the monotone polar and representable closures of monotone operators,
\emph{Journal of Convex Analysis} \textbf{21} (2014), 495--505.

\bibitem{rockafellar1970}
R.~T. Rockafellar,
On the maximality of sums of nonlinear monotone operators,
\emph{Transactions of the American Mathematical Society}
\textbf{149} (1970), 75--88.

\bibitem{yao2010}
L. Yao,
The sum of a maximal monotone operator of type (FPV) and a maximal
monotone operator with full domain is maximal monotone,
\emph{Nonlinear Analysis: Theory, Methods \& Applications}
\textbf{74} (2011), 6144--6152.
doi:10.1016/j.na.2011.05.093.

\bibitem{borweinyao2012}
J.~M. Borwein and L. Yao,
Maximality of the sum of a maximally monotone linear relation and a
maximally monotone operator,
\emph{Set-Valued and Variational Analysis} \textbf{21} (2013), 603--616.
doi:10.1007/s11228-013-0259-y.

\bibitem{voisei2010}
M.~D. Voisei,
A sum theorem for (FPV) operators and normal cones,
\emph{Journal of Mathematical Analysis and Applications}
\textbf{371} (2010), 661--664.
doi:10.1016/j.jmaa.2010.05.064.

\bibitem{marquessvaiter2012}
M. Marques Alves and B.~F. Svaiter,
A new qualification condition for the maximality of the sum of maximal
monotone operators in general Banach spaces,
\emph{Journal of Convex Analysis} \textbf{19} (2012), 575--589.

\bibitem{gossez1971}
J.-P. Gossez,
Op\'{e}rateurs monotones non lin\'{e}aires dans les espaces de Banach non r\'{e}flexifs,
\emph{Journal of Mathematical Analysis and Applications}
\textbf{34} (1971), 371--395.
doi:10.1016/0022-247X(71)90119-3.

\bibitem{marquessvaiter2010}
M. Marques Alves and B.~F. Svaiter,
On Gossez type (D) maximal monotone operators,
\emph{Journal of Convex Analysis} \textbf{17} (2010), 1077--1088.

\bibitem{bauschkefp2012}
H.~H. Bauschke, J.~M. Borwein, X. Wang, and L. Yao,
Every maximally monotone operator of Fitzpatrick--Phelps type is actually
of dense type,
\emph{Optimization Letters} \textbf{6} (2012), 1875--1881.
doi:10.1007/s11590-011-0383-2.

\bibitem{buenosvaiter2011}
O. Bueno and B.~F. Svaiter,
A non-type (D) operator in $c_0$,
\emph{Mathematical Programming} \textbf{139} (2013), 81--88.
doi:10.1007/s10107-013-0661-0.
Lemma numbering refers to the author version, arXiv:1103.2349v1.

\bibitem{bauschkeetal2011}
H.~H. Bauschke, J.~M. Borwein, X. Wang, and L. Yao,
Construction of pathological maximally monotone operators on
non-reflexive Banach spaces,
\emph{Set-Valued and Variational Analysis} \textbf{20} (2012), 387--415.
doi:10.1007/s11228-012-0209-0.
Example numbering refers to the authors' August~6, 2011 version.

\bibitem{eberhardwenczel2024}
A. Eberhard and R. Wenczel,
Representative functions, variational convergence and almost convexity,
\emph{Set-Valued and Variational Analysis} \textbf{32} (2024), article~4.
doi:10.1007/s11228-024-00708-4.

\bibitem{botcsetnek2008}
R.~I. Bo\c{t} and E.~R. Csetnek,
Weak regularity conditions for maximal monotonicity in separable
Asplund spaces,
Chemnitz Faculty of Mathematics, Preprint~3, April~1, 2008.

\bibitem{verona2008}
A. Verona and M. E. Verona,
Regular maximal monotone multifunctions and enlargements,
\emph{Journal of Convex Analysis} \textbf{16} (2009), 1003--1009.

\bibitem{rockafellar1970subdiff}
R. T. Rockafellar,
On the maximal monotonicity of subdifferential mappings,
\emph{Pacific Journal of Mathematics} \textbf{33} (1970), 209--216.

\bibitem{botbuenosimons2016}
R. I. Bo\c{t}, O. Bueno and S. Simons,
The Rockafellar conjecture and type (FPV),
\emph{Set-Valued and Variational Analysis} \textbf{24} (2016), 381--385.
doi:10.1007/s11228-016-0384-5.

\end{thebibliography}

\end{document}
